%Môn: Địa Lí
\begin{ex}
    Dựa vào chỉ số GNI/người năm 2020, quốc gia nào sau đây được xếp vào nhóm nước phát triển?
    \choice
    {Việt Nam}
    {\True Hoa Kỳ}
    {Bra-xin}
    {Nam Phi}

    \loigiai{Theo bảng 1.1 trang 8, Hoa Kỳ có GNI/người là 64.140 USD, thuộc nhóm thu nhập cao và là nước phát triển.}
\end{ex}
%Môn: Địa Lí
\begin{ex}
    Trong cơ cấu GDP của các nước đang phát triển, ngành kinh tế nào thường chiếm tỉ trọng cao nhất nhưng đang có xu hướng giảm?
    \choice
    {Dịch vụ}
    {Công nghiệp}
    {\True Nông nghiệp}
    {Xây dựng}

    \loigiai{Ở các nước đang phát triển, tỉ trọng ngành nông nghiệp, lâm nghiệp và thủy sản thường còn cao hơn so với các nước phát triển nhưng đang có xu hướng giảm dần trong quá trình công nghiệp hóa.}
\end{ex}
%Môn: Địa Lí
\begin{ex}
    Cho bảng số liệu về chỉ số HDI của một số quốc gia năm 2020:
    \choiceTF
    {\True Đức là quốc gia có chỉ số HDI cao nhất trong bảng (0,944).}
    {Việt Nam thuộc nhóm có chỉ số HDI rất cao.}
    {\True Hoa Kỳ và Nhật Bản đều có chỉ số HDI trên 0,800.}
    {Bra-xin có chỉ số HDI thấp hơn Việt Nam.}

    \loigiai{
        \begin{itemchoice}
        \itemch theo bảng 1.1 trang 8.
        \itemch Việt Nam có HDI là 0,710, thuộc nhóm cao nhưng chưa phải rất cao.
        \itemch Hoa Kỳ (0,920) và Nhật Bản (0,923) đều > 0,800.
        \itemch Bra-xin (0,758) cao hơn Việt Nam (0,710).
        \end{itemchoice}
    }
\end{ex}
%Môn: Hóa Học
\begin{ex}
    Hằng số cân bằng $K_c$ của một phản ứng thuận nghịch phụ thuộc vào yếu tố nào sau đây?
    \choice
    {Nồng độ}
    {Áp suất}
    {Xúc tác}
    {\True Nhiệt độ}

    \loigiai{Hằng số cân bằng $K_c$ của phản ứng thuận nghịch chỉ phụ thuộc vào bản chất của phản ứng và nhiệt độ.}
\end{ex}
%Môn: Hóa Học
\begin{ex}
    Trong dung dịch, chất nào sau đây là chất điện li mạnh?
    \choice
    {$CH_3COOH$}
    {$H_2O$}
    {\True $NaOH$}
    {$Mg(OH)_2$}

    \loigiai{$NaOH$ là một base mạnh, phân li hoàn toàn trong nước nên là chất điện li mạnh.}
\end{ex}
%Môn: Hóa Học
\begin{ex}
    Tính pH của dung dịch $NaOH$ 0,01 M ở $25^\circ C$.
    \shortans{12}
    \loigiai{$[OH^-] = 0,01 = 10^{-2} M \Rightarrow pOH = 2 \Rightarrow pH = 14 - 2 = 12$.}
\end{ex}
%Môn: Lịch Sử
\begin{ex}
    Sự kiện nào sau đây được coi là mở đầu cho Cách mạng tư sản Pháp năm 1789?
    \choice
    {Vua Lu-i XVI triệu tập Hội nghị ba đẳng cấp}
    {\True Quần chúng nhân dân tấn công ngục Ba-xti}
    {Quốc hội lập hiến thông qua Tuyên ngôn Nhân quyền và Dân quyền}
    {Phái Gia-cô-banh lên nắm chính quyền}

    \loigiai{Ngày 14/7/1789, quần chúng nhân dân Pa-ri tấn công ngục Ba-xti, biểu tượng của chế độ phong kiến chuyên chế, mở đầu cho cuộc đại cách mạng.}
\end{ex}
%Môn: Lịch Sử
\begin{ex}
    Quốc gia nào ở Đông Nam Á không trở thành thuộc địa của thực dân phương Tây?
    \choice
    {Việt Nam}
    {In-đô-nê-xi-a}
    {\True Xiêm (Thái Lan)}
    {Phi-líp-pin}

    \loigiai{Xiêm là quốc gia duy nhất ở Đông Nam Á giữ được độc lập tương đối nhờ chính sách ngoại giao khôn khéo và cải cách của vua Ra-ma V.}
\end{ex}
%Môn: Lịch Sử
\begin{ex}
    \sochc{2}{
        Dưới thời nhà Nguyễn, các đội Hoàng Sa và Bắc Hải được thành lập để thực thi chủ quyền của Việt Nam tại quần đảo Hoàng Sa và Trường Sa.
    }
    \begin{chc}
        Nhiệm vụ chính của đội Hoàng Sa là gì?
        \loigiai{Nhiệm vụ chính là thu lượm hàng hóa của các tàu bị đắm, đo đạc thủy trình, vẽ bản đồ và khai thác sản vật tại quần đảo Hoàng Sa.}
    \end{chc}
    \begin{chc}
        Đội Bắc Hải chịu sự quản lí trực tiếp của đội nào?
        \loigiai{Đội Bắc Hải chịu sự quản lí của đội Hoàng Sa.}
    \end{chc}
\end{ex}
%Môn: Vật Lí
\begin{ex}
    Một vật dao động điều hòa theo phương trình $x = 5\cos(4\pi t + \pi/3)$ cm. Biên độ của dao động là
    \choice
    {$4\pi$ cm}
    {$\pi/3$ cm}
    {\True 5 cm}
    {10 cm}

    \loigiai{Từ phương trình $x = A\cos(\omega t + \varphi)$, ta có $A = 5$ cm.}
\end{ex}
%Môn: Vật Lí
\begin{ex}
    Hiện tượng cộng hưởng xảy ra khi nào?
    \choice
    {Tần số của ngoại lực rất lớn}
    {Tần số của ngoại lực rất nhỏ}
    {\True Tần số của ngoại lực bằng tần số dao động riêng của hệ}
    {Biên độ của ngoại lực bằng biên độ dao động riêng}

    \loigiai{Hiện tượng biên độ dao động cưỡng bức tăng đến giá trị cực đại khi tần số của ngoại lực bằng tần số riêng của hệ gọi là hiện tượng cộng hưởng.}
\end{ex}
%Môn: Vật Lí
\begin{ex}
    Tại vị trí cân bằng, vận tốc của vật dao động điều hòa có độ lớn bằng bao nhiêu?
    \loigiai{Ở vị trí cân bằng ($x=0$), động năng cực đại và vận tốc có độ lớn cực đại $v = \omega A$.}
\end{ex}
%Môn: Toán
\begin{ex}
    Tính đạo hàm của hàm số $y = x^3 + 2x$.
    \loigiai{Áp dụng quy tắc tính đạo hàm: $(x^n)' = nx^{n-1}$ và $(kx)' = k$.}
\end{ex}
%Môn: Toán
\begin{ex}
    Cho cấp số cộng $(u_n)$ có $u_1 = 3$ và công sai $d = 2$. Tìm số hạng thứ 5 của dãy số.
    \choice
    {7}
    {9}
    {\True 11}
    {13}

    \loigiai{$u_5 = u_1 + 4d = 3 + 4 \cdot 2 = 11$.}
\end{ex}
%Môn: Toán
\begin{ex}
    Trong không gian, nếu đường thẳng $a$ vuông góc với hai đường thẳng cắt nhau nằm trong mặt phẳng $(P)$ thì đường thẳng $a$ vuông góc với mặt phẳng $(P)$.
    \choiceTF
    {\True Đây là điều kiện để đường thẳng vuông góc với mặt phẳng.}
    {Đường thẳng $a$ chỉ cần vuông góc với một đường thẳng trong $(P)$.}
    {Đường thẳng $a$ phải vuông góc với mọi đường thẳng trong không gian.}
    {\True Nếu $a \perp (P)$ thì $a$ vuông góc với mọi đường thẳng nằm trong $(P)$.}

    \loigiai{
        \begin{itemchoice}
        \itemch Đúng theo định lí.
        \itemch phải là hai đường thẳng cắt nhau.
        \itemch chỉ cần trong mặt phẳng đó.
        \itemch Đúng theo tính chất.
        \end{itemchoice}
    }
\end{ex}
%Môn: Giáo Dục Kinh Tế và Pháp Luật
\begin{ex}
    Tình trạng mức giá chung của nền kinh tế tăng lên liên tục trong một khoảng thời gian nhất định được gọi là
    \choice
    {Thất nghiệp}
    {\True Lạm phát}
    {Cạnh tranh}
    {Cung cầu}

    \loigiai{Lạm phát là sự tăng mức giá chung các mặt hàng của nền kinh tế một cách liên tục trong một thời gian nhất định.}
\end{ex}
%Môn: Giáo Dục Kinh Tế và Pháp Luật
\begin{ex}
    Hành vi nào sau đây là biểu hiện của cạnh tranh không lành mạnh?
    \choice
    {Nâng cao chất lượng sản phẩm}
    {Giảm giá thành nhờ cải tiến kỹ thuật}
    {\True Gièm pha đối thủ cạnh tranh bằng cách tung tin sai sự thật}
    {Tăng cường quảng cáo trung thực}

    \loigiai{Gièm pha, nói xấu đối thủ là hành vi vi phạm đạo đức kinh doanh và pháp luật về cạnh tranh.}
\end{ex}
%Môn: Sinh Học
\begin{ex}
    Sắc tố nào trực tiếp tham gia chuyển hóa năng lượng ánh sáng thành năng lượng hóa học trong ATP và NADPH?
    \choice
    {Carotenoid}
    {Diệp lục b}
    {\True Diệp lục a}
    {Xanthophyll}

    \loigiai{Chỉ có diệp lục a ở trung tâm phản ứng mới có khả năng chuyển hóa quang năng thành hóa năng.}
\end{ex}
%Môn: Sinh Học
\begin{ex}
    Sản phẩm của pha sáng trong quang hợp gồm những gì?
    \loigiai{Pha sáng sử dụng nước và ánh sáng để tạo ra năng lượng dưới dạng ATP, NADPH và giải phóng khí oxy.}
\end{ex}
%Môn: Tiếng Anh
\begin{ex}
    Which of the following is considered a healthy habit?
    \choice
    {Staying up late}
    {Eating too much fast food}
    {\True Doing regular exercise}
    {Spending too much screen time}

    \loigiai{Unit 1 highlights regular exercise and a balanced diet as keys to a long and healthy life.}
\end{ex}
%Môn: Tiếng Anh
\begin{ex}
    ASEAN stands for the Association of Southeast Asian ______.
    \choice
    {Nations}
    {\True Nations}
    {Neighborhoods}
    {Networks}

    \loigiai{ASEAN: Association of Southeast Asian Nations.}
\end{ex}
%Môn: Địa Lí
\begin{ex}
    \immini[thm]{
        Dựa vào bảng 7.1, hãy xác định quốc gia có GDP lớn nhất khu vực Mỹ La-tinh năm 2020.
    }{
        [Bảng 7.1]
    }
    \choice
    {Ác-hen-ti-na}
    {Chi-lê}
    {Mê-hi-cô}
    {\True Bra-xin}

    \loigiai{Theo bảng 7.1 trang 30, Bra-xin có GDP cao nhất với 1448,7 tỉ USD.}
\end{ex}
%Môn: Toán
\begin{ex}
    Tính giới hạn: $\lim_{n \to \infty} \frac{1}{n}$.
    \shortans{0}
    \loigiai{Khi n tiến tới vô cùng, nghịch đảo của n tiến về 0.}
\end{ex}
%Môn: Vật Lí
\begin{ex}
    \sochc{2}{
        Sóng cơ truyền trên một sợi dây đàn hồi.
    }
    \begin{chc}
        Khoảng cách giữa hai đỉnh sóng liên tiếp là gì?
        \loigiai{Là bước sóng ($\lambda$).}
    
    \end{chc}
    \begin{chc}
        Tốc độ truyền sóng phụ thuộc vào yếu tố nào?
        \loigiai{Phụ thuộc vào bản chất của môi trường truyền sóng.}
    \end{chc}
\end{ex}
%Môn: Hóa Học
\begin{ex}
    Dẫn xuất halogen của hydrocarbon là hợp chất có nguyên tử nào thay thế nguyên tử hydrogen?
    \choice
    {Oxygen}
    {Nitrogen}
    {\True Halogen}
    {Sulfur}

    \loigiai{Khi thay thế nguyên tử hydrogen trong phân tử hydrocarbon bằng nguyên tử halogen, ta được dẫn xuất halogen.}
\end{ex}
%Môn: Giáo Dục Kinh Tế và Pháp Luật
\begin{ex}
    Quyền nào sau đây thuộc nhóm quyền tự do cơ bản của công dân?
    \choice
    {Quyền bầu cử}
    {Quyền ứng cử}
    {\True Quyền bất khả xâm phạm về thân thể}
    {Quyền tham gia quản lý nhà nước}

    \loigiai{Quyền bất khả xâm phạm về thân thể là một trong những quyền tự do cá nhân quan trọng nhất.}
\end{ex}
%Môn: Lịch Sử
\begin{ex}
    Cuộc Duy tân Minh Trị ở Nhật Bản diễn ra vào thời gian nào?
    \choice
    {1858}
    {\True 1868}
    {1888}
    {1911}

    \loigiai{Cuộc Duy tân Minh Trị bắt đầu từ năm 1868, đưa Nhật Bản phát triển theo con đường tư bản chủ nghĩa.}
\end{ex}
%Môn: Địa Lí
\begin{ex}
    Liên minh châu Âu (EU) được thành lập chính thức vào năm nào qua Hiệp ước Ma-xtrích?
    \choice
    {1957}
    {1967}
    {\True 1993}
    {2004}

    \loigiai{Hiệp ước Ma-xtrích có hiệu lực từ ngày 1-11-1993, đánh dấu sự ra đời chính thức của EU.}
\end{ex}
%Môn: Hóa Học
\begin{ex}
    Chất nào sau đây là thành phần chính của khí thiên nhiên?
    \choice
    {Ethane}
    {Propane}
    {\True Methane}
    {Butane}

    \loigiai{Methane ($CH_4$) chiếm thành phần chủ yếu trong khí thiên nhiên.}
\end{ex}
%Môn: Vật Lí
\begin{ex}
    Đơn vị của cường độ điện trường là
    \choice
    {Vôn (V)}
    {Ampe (A)}
    {\True Vôn trên mét (V/m)}
    {Cu-lông (C)}

    \loigiai{Cường độ điện trường được đo bằng V/m.}
\end{ex}
%Môn: Toán
\begin{ex}
    Giải phương trình: $2^x = 8$.
    \shortans{3}
    \loigiai{$2^x = 2^3 \Rightarrow x = 3$.}
\end{ex}
%Môn: Sinh Học
\begin{ex}
    Hô hấp ở thực vật diễn ra chủ yếu ở bào quan nào?
    \choice
    {Lục lạp}
    {Không bào}
    {\True Ti thể}
    {Nhân tế bào}

    \loigiai{Ti thể là nơi diễn ra quá trình phân giải chất hữu cơ để giải phóng năng lượng.}
\end{ex}
%Môn: Giáo Dục Kinh Tế và Pháp Luật
\begin{ex}
    Khi lượng cung lớn hơn lượng cầu trên thị trường thì giá cả mặt hàng đó thường có xu hướng
    \choice
    {Tăng}
    {\True Giảm}
    {Không đổi}
    {Biến động mạnh}

    \loigiai{Theo quy luật cung - cầu, khi cung > cầu, giá cả thường giảm xuống để kích cầu.}
\end{ex}
%Môn: Địa Lí
\begin{ex}
    Bra-xin là quốc gia thuộc khu vực nào?
    \choice
    {Đông Nam Á}
    {Tây Nam Á}
    {\True Mỹ La-tinh}
    {Liên minh châu Âu}

    \loigiai{Bra-xin là quốc gia lớn nhất ở khu vực Mỹ La-tinh.}
\end{ex}
%Môn: Tiếng Anh
\begin{ex}
    The generation gap refers to the difference in ______ between young and old people.
    \choice
    {Height}
    {\True Attitudes and behaviors}
    {Physical strength}
    {Hair color}

    \loigiai{Unit 2 defines the generation gap as differences in beliefs, attitudes, and behaviors.}
\end{ex}
%Môn: Toán
\begin{ex}
    Trong các hàm số sau, hàm số nào là hàm số chẵn?
    \choice
    {$y = \sin x$}
    {\True $y = \cos x$}
    {$y = \tan x$}
    {$y = \cot x$}

    \loigiai{Hàm số $y = \cos x$ thỏa mãn $\cos(-x) = \cos x$ nên là hàm số chẵn.}
\end{ex}
%Môn: Hóa Học
\begin{ex}
    Nguyên tố nào là thành phần không thể thiếu trong mọi hợp chất hữu cơ?
    \choice
    {Oxygen}
    {Hydrogen}
    {Nitrogen}
    {\True Carbon}

    \loigiai{Hợp chất hữu cơ là hợp chất của carbon (trừ một số ngoại lệ như $CO, CO_2$, muối carbonate...)}
\end{ex}
%Môn: Vật Lí
\begin{ex}
    Trong hiện tượng giao thoa sóng mặt nước, những điểm cực đại giao thoa là những điểm mà hai sóng tới tại đó
    \choice
    {Ngược pha}
    {\True Cùng pha}
    {Lệch pha $\pi/2$}
    {Lệch pha $\pi/4$}

    \loigiai{Hai sóng cùng pha tăng cường lẫn nhau tạo nên cực đại.}
\end{ex}
%Môn: Lịch Sử
\begin{ex}
    Cuộc khởi nghĩa Lam Sơn (1418 - 1427) chống lại quân xâm lược nào?
    \choice
    {Nhà Tống}
    {Nhà Nguyên}
    {\True Nhà Minh}
    {Nhà Thanh}

    \loigiai{Cuộc khởi nghĩa Lam Sơn do Lê Lợi lãnh đạo nhằm đánh đuổi quân Minh.}
\end{ex}
%Môn: Sinh Học
\begin{ex}
    Quá trình truyền tin qua synapse hóa học diễn ra theo chiều nào?
    \choice
    {Từ màng sau sang màng trước}
    {\True Từ màng trước sang màng sau}
    {Hai chiều}
    {Ngẫu nhiên}

    \loigiai{Thông tin chỉ được truyền một chiều nhờ các túi chứa chất trung gian hóa học ở màng trước.}
\end{ex}
%Môn: Giáo Dục Kinh Tế và Pháp Luật
\begin{ex}
    Thị trường lao động là nơi thực hiện các quan hệ giữa người lao động và
    \choice
    {Người tiêu dùng}
    {Cửa hàng bán lẻ}
    {\True Người sử dụng lao động}
    {Cơ quan nhà nước}

    \loigiai{Thị trường lao động là nơi diễn ra sự thỏa thuận giữa người bán sức lao động và người mua sức lao động.}
\end{ex}
%Môn: Địa Lí
\begin{ex}
    Kênh đào nào nối liền Địa Trung Hải với Biển Đỏ?
    \choice
    {Kênh đào Pa-na-ma}
    {\True Kênh đào Xuy-ê}
    {Kênh đào Ki-en}
    {Kênh đào Cor-inth}

    \loigiai{Kênh đào Xuy-ê nằm ở Ai Cập, nối liền hai vùng biển này, có ý nghĩa quan trọng trong hàng hải.}
\end{ex}
%Môn: Tiếng Anh
\begin{ex}
    Which energy source is renewable?
    \choice
    {Coal}
    {Oil}
    {\True Solar power}
    {Natural gas}

    \loigiai{Unit 5 mentions solar, wind, and hydro as renewable energy sources.}
\end{ex}
%Môn: Toán
\begin{ex}
    Cho đường thẳng $d$ song song với mặt phẳng $(P)$. Có bao nhiêu mặt phẳng chứa $d$ và song song với $(P)$?
    \choice
    {0}
    {\True 1}
    {2}
    {Vô số}

    \loigiai{Chỉ có duy nhất một mặt phẳng thỏa mãn điều kiện này.}
\end{ex}
%Môn: Hóa Học
\begin{ex}
    Phản ứng nào sau đây được dùng để sản xuất ammonia trong công nghiệp?
    \choice
    {$N_2 + O_2$}
    {\True $N_2 + H_2$}
    {$NH_4Cl + NaOH$}
    {$NH_3 + HCl$}

    \loigiai{ Ammonia được tổng hợp theo quá trình Haber-Bosch từ nitơ và hiđrô.}
\end{ex}
%Môn: Vật Lí
\begin{ex}
    Âm thanh không truyền được trong môi trường nào?
    \choice
    {Rắn}
    {Lỏng}
    {Khí}
    {\True Chân không}

    \loigiai{Sóng âm cần môi trường vật chất để lan truyền.}
\end{ex}
%Môn: Lịch Sử
\begin{ex}
    Tác phẩm nào của Nguyễn Ái Quốc đã tố cáo tội ác của thực dân Pháp?
    \loigiai{Tác phẩm này được xuất bản lần đầu tại Pháp năm 1925.}
\end{ex}
%Môn: Sinh Học
\begin{ex}
    Hiện tượng lá cây trinh nữ cụp lại khi chạm vào là ví dụ về
    \choice
    {Hướng sáng}
    {Hướng nước}
    {\True Ứng động}
    {Hướng hóa}

    \loigiai{Đây là ứng động không sinh trưởng (ứng động tiếp xúc).}
\end{ex}
%Môn: Giáo Dục Kinh Tế và Pháp Luật
\begin{ex}
    \sochc{2}{
        Hành vi trốn thuế của doanh nghiệp gây thiệt hại cho ngân sách nhà nước và xã hội.
    }
    \begin{chc}
        Hành vi trốn thuế có vi phạm đạo đức kinh doanh không?
        \loigiai{Có, vì nó vi phạm tính trung thực và trách nhiệm với cộng đồng.}
    \end{chc}
    \begin{chc}
        Cơ quan nào có thẩm quyền xử lý hành vi này?
        \loigiai{Cơ quan Thuế và các cơ quan chức năng có thẩm quyền theo quy định của pháp luật.}
    \end{chc}
\end{ex}
%Môn: Địa Lí
\begin{ex}
    Sông nào là sông dài nhất thế giới chảy qua khu vực phía Bắc châu Phi?
    \loigiai{Sông Nin là con sông quan trọng nhất của khu vực này.}
\end{ex}
%Môn: Tiếng Anh
\begin{ex}
    What does a 'smart city' use to operate efficiently?
    \choice
    {Old buses}
    {\True AI and sensor technologies}
    {Coal power}
    {Manual labor only}

    \loigiai{Unit 3 explains that smart cities use technology like AI and sensors for management.}
\end{ex}
%Môn: Toán
\begin{ex}
    Tính giới hạn của dãy số $u_n = \frac{2n+1}{n+3}$.
    \shortans{2}
    \loigiai{Chia cả tử và mẫu cho n: $\lim \frac{2 + 1/n}{1 + 3/n} = 2$.}
\end{ex}
%Môn: Hóa Học
\begin{ex}
    Phương pháp nào được dùng để tách các chất lỏng không hòa tan vào nhau?
    \choice
    {Chưng cất}
    {\True Chiết}
    {Kết tinh}
    {Sắc ký}

    \loigiai{Phương pháp chiết lỏng - lỏng dùng để tách hai chất lỏng không trộn lẫn.}
\end{ex}
%Môn: Vật Lí
\begin{ex}
    Sóng điện từ là sóng
    \choice
    {Sóng dọc}
    {\True Sóng ngang}
    {Sóng dừng}
    {Sóng âm}

    \loigiai{Sóng điện từ là sóng ngang vì các thành phần điện trường và từ trường vuông góc với phương truyền sóng.}
\end{ex}
%Môn: Lịch Sử
\begin{ex}
    Vua nào thực hiện cuộc cải cách hành chính chia cả nước thành 30 tỉnh và 1 phủ Thừa Thiên vào năm 1831 - 1832?
    \choice
    {Gia Long}
    {\True Minh Mạng}
    {Thiệu Trị}
    {Tự Đức}

    \loigiai{Vua Minh Mạng thực hiện cải cách lớn về hành chính trong nửa đầu thế kỉ XIX.}
\end{ex}
%Môn: Sinh Học
\begin{ex}
    Hormone thực vật nào kích thích sự chín của quả?
    \choice
    {Auxin}
    {Gibberellin}
    {Cytokinin}
    {\True Ethylene}

    \loigiai{Ethylene là hormone dạng khí thúc đẩy quá trình chín của quả.}
\end{ex}
%Môn: Giáo Dục Kinh Tế và Pháp Luật
\begin{ex}
    Người đủ bao nhiêu tuổi trở lên thì có quyền bầu cử đại biểu Quốc hội?
    \choice
    {16 tuổi}
    {\True 18 tuổi}
    {21 tuổi}
    {25 tuổi}

    \loigiai{Theo Hiến pháp, công dân đủ 18 tuổi trở lên có quyền bầu cử.}
\end{ex}
%Môn: Địa Lí
\begin{ex}
    Arene là loại hydrocarbon có chứa cấu trúc đặc trưng nào?
    \choice
    {Mạch hở}
    {Liên kết ba}
    {\True Vòng benzene}
    {Liên kết đôi}

    \loigiai{Arene (hydrocarbon thơm) là những hydrocarbon có chứa vòng benzene trong phân tử.}
\end{ex}
%Môn: Toán
\begin{ex}
    Cho hai biến cố độc lập $A$ và $B$ với $P(A) = 0,5$ và $P(B) = 0,2$. Tính $P(AB)$.
    \shortans{0,1}
    \loigiai{Vì A, B độc lập nên $P(AB) = P(A) \cdot P(B) = 0,5 \cdot 0,2 = 0,1$.}
\end{ex}
%Môn: Vật Lí
\begin{ex}
    Một điện tích điểm $q$ đặt trong điện trường chịu tác dụng của lực điện $F$. Cường độ điện trường $E$ tại điểm đó được tính theo công thức:
    \choice
    {$E = qF$}
    {\True $E = F/q$}
    {$E = q/F$}
    {$E = F^2/q$}

    \loigiai{Định nghĩa cường độ điện trường $E = F/q$.}
\end{ex}
%Môn: Hóa Học
\begin{ex}
    Công thức của Sulfuric acid là
    \loigiai{Axit sunfuric có công thức hóa học là $H_2SO_4$.}
\end{ex}
%Môn: Sinh Học
\begin{ex}
    Loại mạch nào dẫn nước và muối khoáng từ rễ lên lá?
    \choice
    {Mạch rây}
    {\True Mạch gỗ}
    {Mạch rây và mạch gỗ}
    {Biểu bì}

    \loigiai{Mạch gỗ (xylem) vận chuyển dòng mạch gỗ gồm nước và khoáng.}
\end{ex}
%Môn: Lịch Sử
\begin{ex}
    Liên bang Cộng hòa xã hội chủ nghĩa Xô viết (Liên Xô) được thành lập vào năm nào?
    \choice
    {1917}
    {1919}
    {1921}
    {\True 1922}

    \loigiai{Liên Xô thành lập vào cuối năm 1922.}
\end{ex}
%Môn: Địa Lí
\begin{ex}
    Quốc gia nào có diện tích rộng nhất thế giới?
    \choice
    {Canada}
    {Hoa Kỳ}
    {Trung Quốc}
    {\True Liên bang Nga}

    \loigiai{Liên bang Nga có diện tích khoảng 17 triệu km2, rộng nhất thế giới.}
\end{ex}
%Môn: Tiếng Anh
\begin{ex}
    People who were born between the early 1980s and late 1990s are called ______.
    \choice
    {Gen Xers}
    {\True Millennials}
    {Gen Zers}
    {Baby Boomers}

    \loigiai{Unit 2 defines Generation Y as Millennials.}
\end{ex}
%Môn: Toán
\begin{ex}
    Cho hình chóp $S.ABCD$ có đáy là hình vuông. Đường thẳng $SA$ vuông góc với mặt phẳng đáy $(ABCD)$. Khẳng định nào sau đây đúng?
    \choice
    {$SA \perp BC$}
    {$SA \perp CD$}
    {$SA \perp BD$}
    {\True Cả A, B, C đều đúng}

    \loigiai{Vì $SA \perp (ABCD)$ nên $SA$ vuông góc với mọi đường thẳng nằm trong mặt phẳng $(ABCD)$.}
\end{ex}
%Môn: Hóa Học
\begin{ex}
    \immini[thm]{
        Cho mô hình phân tử sau. Đây là chất nào?
    }{
        [Mô hình CH4]}
    \choice
    {Ethane}
    {Ethylene}
    {\True Methane}
    {Acetylene}

    \loigiai{Mô hình gồm 1 nguyên tử Carbon trung tâm liên kết với 4 nguyên tử Hydrogen là Methane.}
\end{ex}
%Môn: Vật Lí
\begin{ex}
    Tụ điện là linh kiện dùng để
    \choice
    {Tăng điện áp}
    {Đốt nóng}
    {\True Tích trữ điện tích và năng lượng điện}
    {Cản trở dòng điện hoàn toàn}

    \loigiai{Tụ điện có khả năng tích điện và phóng điện.}
\end{ex}
%Môn: Lịch Sử
\begin{ex}
    Cuộc cải cách của Hồ Quý Ly diễn ra vào khoảng thời gian nào?
    \choice
    {Đầu thế kỉ XIV}
    {\True Cuối thế kỉ XIV - đầu thế kỉ XV}
    {Giữa thế kỉ XV}
    {Cuối thế kỉ XV}

    \loigiai{Cuộc cải cách diễn ra vào cuối thời Trần và trong thời Hồ (1400-1407).}
\end{ex}
%Môn: Sinh Học
\begin{ex}
    Thú có vú hô hấp bằng bộ phận nào?
    \choice
    {Da}
    {Mang}
    {Hệ thống ống khí}
    {\True Phổi}

    \loigiai{Thú là nhóm động vật hô hấp bằng phổi có hiệu quả nhất.}
\end{ex}
%Môn: Địa Lí
\begin{ex}
    Sự kiện nào đánh dấu sự chấm dứt của chiến tranh lạnh?
    \loigiai{Sự sụp đổ của hệ thống xã hội chủ nghĩa ở Đông Âu và sự tan rã của Liên Xô (1989-1991).}
\end{ex}
%Môn: Giáo Dục Kinh Tế và Pháp Luật
\begin{ex}
    Thị trường gồm những yếu tố cơ bản nào?
    \choice
    {Chỉ có người mua và người bán}
    {Chỉ có hàng hóa và tiền tệ}
    {\True Người mua, người bán, hàng hóa, tiền tệ và giá cả}
    {Chỉ có giá cả và chất lượng}

    \loigiai{Thị trường là tổng thể các mối quan hệ kinh tế giữa các chủ thể.}
\end{ex}
%Môn: Toán
\begin{ex}
    Tính đạo hàm của hàm số $y = \sin x$.
    \loigiai{Đây là công thức đạo hàm cơ bản của hàm số lượng giác.}
\end{ex}
%Môn: Hóa Học
\begin{ex}
    Nitrogen có công thức hóa học là
    \loigiai{Nitơ điôxít là $NO_2$.}
\end{ex}
%Môn: Vật Lí
\begin{ex}
    Sóng dừng được hình thành do sự kết hợp của
    \choice
    {Hai sóng cùng tần số}
    {Hai sóng cùng biên độ}
    {\True Sóng tới và sóng phản xạ cùng tần số}
    {Hai sóng siêu âm}

    \loigiai{Sóng dừng là kết quả của sự giao thoa giữa sóng tới và sóng phản xạ.}
\end{ex}
%Môn: Lịch Sử
\begin{ex}
    Nguyên nhân chủ yếu dẫn đến sự thất bại của các cuộc kháng chiến chống ngoại xâm là do
    \choice
    {Vũ khí thô sơ}
    {Thiếu nhân tài}
    {\True Sai lầm về đường lối và sự thiếu đoàn kết}
    {Quân giặc quá đông}

    \loigiai{Lịch sử cho thấy sự đồng lòng và đường lối đúng đắn là quyết định nhất.}
\end{ex}
%Môn: Sinh Học
\begin{ex}
    Pheromone là gì?
    \loigiai{Pheromone là chất hóa học do động vật sản sinh và giải phóng vào môi trường để gây ra các đáp ứng khác nhau ở các cá thể cùng loài.}
\end{ex}
%Môn: Giáo Dục Kinh Tế và Pháp Luật
\begin{ex}
    Công dân có nghĩa vụ gì đối với nhà nước?
    \choice
    {Học tập}
    {Lao động}
    {\True Nộp thuế theo quy định}
    {Vui chơi giải trí}

    \loigiai{Nộp thuế là nghĩa vụ bắt buộc để nhà nước duy trì hoạt động và phúc lợi.}
\end{ex}
%Môn: Địa Lí
\begin{ex}
    Tổ chức WTO có trụ sở tại quốc gia nào?
    \choice
    {Hoa Kỳ}
    {Pháp}
    {\True Thụy Sĩ}
    {Anh}

    \loigiai{WTO có trụ sở tại Giơ-ne-vơ, Thụy Sĩ.}
\end{ex}
%Môn: Tiếng Anh
\begin{ex}
    Global warming leads to rising ______.
    \choice
    {Mountains}
    {\True Sea levels}
    {Buildings}
    {Populations}

    \loigiai{Unit 5 explains that as temperatures rise, ice caps melt, raising sea levels.}
\end{ex}
%Môn: Toán
\begin{ex}
    Trong không gian, cho hai đường thẳng phân biệt cùng vuông góc với một mặt phẳng. Khi đó hai đường thẳng đó
    \choice
    {Cắt nhau}
    {Chéo nhau}
    {\True Song song với nhau}
    {Trùng nhau}

    \loigiai{Đây là tính chất liên hệ giữa quan hệ song song và quan hệ vuông góc.}
\end{ex}
%Môn: Hóa Học
\begin{ex}
    Rượu bia có chứa chất nào gây hại cho gan và thần kinh?
    \loigiai{Ethanol ($C_2H_5OH$) là thành phần chính của đồ uống có cồn.}
\end{ex}
%Môn: Vật Lí
\begin{ex}
    Gia tốc của vật dao động điều hòa biến thiên như thế nào so với li độ?
    \choice
    {Cùng pha}
    {\True Ngược pha}
    {Sớm pha $\pi/2$}
    {Trễ pha $\pi/2$}

    \loigiai{Công thức $a = -\omega^2 x$.}
\end{ex}
%Môn: Lịch Sử
\begin{ex}
    Ai là người chỉ huy chiến thắng Bạch Đằng năm 938?
    \loigiai{Chiến thắng Bạch Đằng năm 938 chấm dứt hơn 1000 năm Bắc thuộc.}
\end{ex}
%Môn: Sinh Học
\begin{ex}
    Đặc điểm nào chỉ có ở sinh sản hữu tính?
    \choice
    {Tạo ra cá thể mới giống hệt mẹ}
    {\True Có sự kết hợp giữa giao tử đực và giao tử cái}
    {Không cần sự tham gia của tế bào sinh dục}
    {Diễn ra rất nhanh}

    \loigiai{Sinh sản hữu tính đặc trưng bởi sự thụ tinh tạo hợp tử.}
\end{ex}
%Môn: Giáo Dục Kinh Tế và Pháp Luật
\begin{ex}
    Thất nghiệp xảy ra khi người lao động có khả năng lao động nhưng không
    \choice
    {Muốn làm việc}
    {\True Tìm được việc làm}
    {Có sức khỏe}
    {Đến tuổi lao động}

    \loigiai{Thất nghiệp là tình trạng người lao động mong muốn có việc làm nhưng chưa tìm được.}
\end{ex}
%Môn: Địa Lí
\begin{ex}
    Nhật Bản là một quốc gia quần đảo nằm ở đại dương nào?
    \choice
    {Ấn Độ Dương}
    {Đại Tây Dương}
    {\True Thái Bình Dương}
    {Bắc Băng Dương}

    \loigiai{Nhật Bản nằm ở phía đông châu Á, bao quanh là Thái Bình Dương.}
\end{ex}
%Môn: Toán
\begin{ex}
    Tính giới hạn của hàm số $f(x) = x^2 - 1$ khi $x \to 2$.
    \shortans{3}
    \loigiai{$2^2 - 1 = 3$.}
\end{ex}
%Môn: Hóa Học
\begin{ex}
    Phản ứng của Phenol với nước Bromine tạo thành kết tủa màu gì?
    \choice
    {Vàng}
    {Đỏ}
    {\True Trắng}
    {Xanh}

    \loigiai{Tạo ra 2,4,6-tribromophenol kết tủa trắng.}
\end{ex}
%Môn: Vật Lí
\begin{ex}
    \immini[thm]{
        Dựa vào đồ thị, hãy xác định biên độ dao động.
    }{
        [Đồ thị hình 2.1]}
    \choice
    {0,1 m}
    {\True 0,2 m}
    {0,3 m}
    {0,4 m}

    \loigiai{Đỉnh cao nhất của đồ thị li độ tương ứng với biên độ $A=0,2$ m.}
\end{ex}
%Môn: Lịch Sử
\begin{ex}
    Cách mạng Tân Hợi diễn ra vào năm nào ở Trung Quốc?
    \shortans{1911}
    \loigiai{Cách mạng này lật đổ triều đại Mãn Thanh.}
\end{ex}
%Môn: Sinh Học
\begin{ex}
    Sự trao đổi khí ở sâu bọ được thực hiện qua
    \choice
    {Phổi}
    {Mang}
    {Da}
    {\True Hệ thống ống khí}

    \loigiai{Ống khí thông ra ngoài qua các lỗ thở.}
\end{ex}
%Môn: Giáo Dục Kinh Tế và Pháp Luật
\begin{ex}
    Quyền bất khả xâm phạm về chỗ ở của công dân có nghĩa là không ai được tự ý vào chỗ ở của người khác nếu không được người đó
    \choice
    {Cho phép}
    {Đồng ý}
    {\True Đồng ý hoặc pháp luật cho phép}
    {Nhìn thấy}

    \loigiai{Việc khám xét chỗ ở phải tuân thủ trình tự pháp luật.}
\end{ex}
%Môn: Địa Lí
\begin{ex}
    Khu vực nào giàu tài nguyên dầu mỏ nhất thế giới?
    \choice
    {Mỹ La-tinh}
    {Đông Nam Á}
    {\True Tây Nam Á}
    {Trung Á}

    \loigiai{Tây Nam Á sở hữu hơn 50% trữ lượng dầu mỏ thế giới.}
\end{ex}
%Môn: Tiếng Anh
\begin{ex}
    Millennials are also known as Digital ______.
    \loigiai{Millennials grew up with technology, so they are digital natives.}
\end{ex}
%Môn: Toán
\begin{ex}
    Tổng của cấp số nhân lùi vô hạn với $|q| < 1$ là
    \choice
    {$S = u_1(1-q)$}
    {\True $S = \frac{u_1}{1-q}$}
    {$S = \frac{u_1}{q-1}$}
    {$S = u_1q^n$}

    \loigiai{Đây là công thức tính tổng cấp số nhân lùi vô hạn.}
\end{ex}
%Môn: Hóa Học
\begin{ex}
    Chất nào sau đây có mùi thơm đặc trưng của quả dứa?
    \choice
    {Ethyl acetate}
    {\True Ethyl butyrate}
    {Isoamyl acetate}
    {Methyl salicylate}

    \loigiai{Ethyl butyrate và ethyl propionate tạo nên mùi dứa.}
\end{ex}
%Môn: Vật Lí
\begin{ex}
    Sóng dọc truyền được trong các môi trường nào?
    \choice
    {Chỉ trong chất khí}
    {Chỉ trong chất lỏng}
    {Chỉ trong chất rắn}
    {\True Rắn, lỏng và khí}

    \loigiai{Sóng dọc lan truyền được trong cả 3 môi trường vật chất.}
\end{ex}
%Môn: Lịch Sử
\begin{ex}
    Quần đảo Trường Sa thuộc tỉnh nào của Việt Nam hiện nay?
    \loigiai{Huyện đảo Trường Sa trực thuộc tỉnh Khánh Hòa.}
\end{ex}
%Môn: Sinh Học
\begin{ex}
    Cây C4 có năng suất cao hơn cây C3 vì
    \choiceTF
    {\True Không có hô hấp sáng.}
    {Không có pha tối.}
    {\True Có điểm bù CO2 thấp.}
    {\True Tận dụng được nồng độ CO2 thấp.}

    \loigiai{
        \begin{itemchoice}
        \itemch C4 tránh được hô hấp sáng.
        \itemch Mọi cây xanh đều có pha tối.
        \itemch 
        \itemch 
        \end{itemchoice}
    }
\end{ex}
%Môn: Giáo Dục Kinh Tế và Pháp Luật
\begin{ex}
    Bình đẳng giới là việc nam, nữ có vị trí, vai trò như thế nào?
    \loigiai{Nam, nữ có vị trí, vai trò ngang nhau, được tạo điều kiện và cơ hội phát huy năng lực như nhau.}
\end{ex}
